\section{Implikationen für Data-Science-Projektmanagement}

Die entwickelte RM-Extension liefert für die Praxis des Data-Science-Projektmanagements 
mehrere konkrete Anknüpfungspunkte. Erstens schafft die Verknüpfung von 
phasenspezifischer Risikobewertung mit formalen Entscheidungspunkten eine 
Grundlage für regelkonforme Projektumsetzungen, insbesondere im Kontext des 
EU AI Acts, der für Hochrisikosysteme explizite Anforderungen an Dokumentation, 
Governance und Monitoring stellt~\cite{EUAIAct2024}. Zweitens unterstützt die 
strukturierte Gate-Logik eine klarere Kommunikation zwischen technischen und 
nicht-technischen Stakeholdern, da Risikostatus und Entscheidungen in einem 
einheitlichen, nachvollziehbaren Format dargestellt werden. Drittens ermöglicht 
die Carry-Forward-Logik eine konsequente Risikoverfolgung über den gesamten 
Projektverlauf, ohne dass erhöhte Risiken zwischen den Phasen in Vergessenheit 
geraten.

Für Organisationen, die die RM-Extension einführen möchten, empfiehlt sich ein 
schrittweises Vorgehen: In einem ersten Schritt sollten die relevanten KRI für 
den jeweiligen Anwendungskontext definiert und, soweit möglich, automatisiert 
erfasst werden. Darauf aufbauend können die Gate-Kriterien kalibriert und in 
bestehende Projektroutinen integriert werden. Eine enge Zusammenarbeit zwischen 
Data Scientists, Risikomanagern und Domänenexpert:innen ist dabei eine 
Grundvoraussetzung, um sowohl technische als auch organisatorische Risiken 
adäquat abzubilden.

Die entwickelte RM-Extension stellt dabei nicht den einzigen denkbaren Integrationsansatz dar. 
Alternativ wäre eine losere Kopplung von Risikomanagement und Prozessmodell denkbar, 
bei der Risikobeurteilungen nicht phasengebunden, sondern ereignisgetrieben durchgeführt 
werden -- etwa ausgelöst durch signifikante Änderungen an Daten, Modellen oder Anforderungen. 
Dieser Ansatz hätte den Vorteil einer höheren Flexibilität in stark iterativen Projekten, 
in denen Phasengrenzen ohnehin fließend sind, ginge jedoch auf Kosten 
der Nachvollziehbarkeit und der strukturierten Dokumentation, die gerade regulatorisch gefordert ist. 
Ein weiterer alternativer Ansatz wäre die vollständige Einbettung des Risikomanagements in bestehende MLOps-Pipelines, 
ohne explizite Anbindung an das DASC-PM. Tools wie MLflow oder Evidently AI ermöglichen ein kontinuierliches, 
automatisiertes Monitoring zentraler Risikoindikatoren im laufenden Betrieb und könnten theoretisch 
als alleiniger Steuerungsmechanismus fungieren. Allerdings decken solche Lösungen vorrangig technische 
Modell- und Datenrisiken ab und vernachlässigen organisatorische, ethische und strategische Risikodimensionen, 
die im DASC-PM explizit adressiert werden. Die RM-Extension positioniert sich 
damit bewusst zwischen diesen Extremen: Sie ist stärker strukturiert als ein rein ereignisgetriebener 
Ansatz und breiter aufgestellt als ein reines MLOps-Monitoring, zahlt dafür jedoch den 
Preis eines höheren initialen Kalibrierungsaufwands. Welcher Ansatz im konkreten Projektkontext geeigneter ist, 
hängt maßgeblich von der Projektgröße, dem Regulierungsumfeld und der organisatorischen Reife im Umgang mit Risikomanagement ab.